\scanpage{96}
% L1–2, centered title
\begin{center}Relations entre les fonctions de Weierstrass et de Legendre\end{center}

% L3–5
Supposons $\omega_1$ réel, $\omega_2$ purement imaginaire ($\unsure{\mathfrak I|J}\omega_2>0$), alors les $e_i$ sont réels; prenons $e_2<e_3<e_1$, posons $k=\sqrt{\frac{e_3-e_2}{e_1-e_2}}$, $M=\sqrt{e_1-e_2}$. Alors

% L6–8
\[
\left\{\begin{aligned}
\operatorname{sn}(u\mid k)&=M\frac{\sigma(\frac uM\mid\omega_1,\omega_2)}{\sigma_2(\frac uM\mid\omega_1,\omega_2)}\\
\operatorname{cn}(u\mid k)&=\deleted{\unsure{M|k}}\frac{\sigma_1(\frac uM\mid\omega_1,\omega_2)}{\sigma_2(\frac uM\mid\omega_1,\omega_2)}\\
\operatorname{dn}(u\mid k)&=\unsure{1|/|,}\,\frac{\sigma_3(\frac uM\mid\omega_1,\omega_2)}{\sigma_1(\frac uM\mid\omega_1,\omega_2)}
\end{aligned}\right.
\]

% L9
\[
\unsure{p|\wp}(u\mid\omega_1,\omega_2)=e_2+\frac{M^2}{\operatorname{sn}^2(Mu\mid k)}
\]

% L10
\[
K=M\omega_1,\quad iK'=M\omega_2,\quad K''=M\omega_3
\]

% L11–16
(Pour $k$ donné, pour \unsure{rattacher} $\operatorname{sn}(u\mid k)\ldots$ aux fonctions de Weierstrass, on prend $M$ arbitraire, puis on détermine $\omega_1,\omega_2,\omega_3$ par $K=M\omega_1$, $iK'=M\omega_{\unsure{2|1}}$, $K''=M\omega_3$, et on a les formules ci-dessus; \deleted{\unsure{où|ici}} alors $e_1,e_2,e_3$ sont donnés par

% L17
\[
3e_1=(1+k'^{\,2})M^2\qquad
3e_2=(\unsure{-2+|2+}k'^{\,2})M^2\qquad
3e_3=(1-2k'^{\,2})M^2
\]

% L18–21, foot
On pourra faire par exemple $M=\frac12$, donc prendre des périodes \inserted{communes} $2\omega_1=4K,\unsure{,|{}}$ $2\omega_2=4iK'$; alors $\operatorname{sn},\operatorname{cn},\operatorname{dn}$ sont fonctions elliptiques d'ordre $4$ pour les périodes $2\omega_1,2\omega_2$)

% lower left
% archive sheet number: 95
